% =====================================================================
%  FICHE 11 — THEME 3 — Exercices DIFFICILE — Équilibrer une équation redox
% =====================================================================
\documentclass[11pt,a4paper]{article}
\usepackage[utf8]{inputenc}
\usepackage[T1]{fontenc}
\usepackage{lmodern}
\usepackage[french]{babel}
\usepackage{amsmath,amssymb}
\usepackage[version=4]{mhchem}
\usepackage{siunitx}
\sisetup{output-decimal-marker={,},per-mode=symbol}
\usepackage{xcolor}
\usepackage{array}
\usepackage{booktabs}
\usepackage{enumitem}
\usepackage[most]{tcolorbox}
\usepackage[a4paper,margin=2.1cm,headheight=26pt,headsep=10pt]{geometry}
\usepackage{fancyhdr}
\usepackage[hidelinks]{hyperref}

\definecolor{primary}{RGB}{146,95,160}
\definecolor{primarydark}{RGB}{92,52,104}
\newcommand{\NO}{\ensuremath{\mathrm{N.O.}}}

\newcounter{exo}
\newtcolorbox{exo}[1][]{enhanced,breakable,boxrule=0.9pt,arc=2pt,colback=primary!4,
  colframe=primary,fonttitle=\bfseries,coltitle=white,left=3mm,right=3mm,top=2.2mm,bottom=2.2mm,
  title={Exercice~\refstepcounter{exo}\theexo\ifx\relax#1\relax\relax\else\ \textmd{--- \itshape #1}\fi}}

\pagestyle{fancy}\fancyhf{}
\fancyhead[L]{\small\textcolor{primarydark}{\textbf{Fiche 11 · Thème 3} — Équilibrer une équation redox}}
\fancyhead[R]{\small\textcolor{primarydark}{\textsc{Difficile}}}
\fancyfoot[C]{\small\thepage}
\renewcommand{\headrulewidth}{0.8pt}
\renewcommand{\headrule}{\hbox to\headwidth{\color{primary}\leaders\hrule height \headrulewidth\hfill}}

\begin{document}
\begin{center}
{\LARGE\bfseries\color{primarydark}Thème 3 — Niveau Difficile}\\[2pt]
{\small\itshape Équilibrages complets, du ionique au moléculaire, en milieu acide et basique.}\\[-4pt]
{\color{primary}\rule{\textwidth}{1pt}}
\end{center}

\begin{exo}[du ionique au moléculaire — d'après annales 2021]
Lorsqu'on mélange du permanganate de potassium et de l'iodure de potassium en présence d'acide
chlorhydrique, il se forme du chlorure de manganèse(II) et du diiode :
\[
\ce{KMnO4 + KI + HCl -> MnCl2 + I2 + KCl + H2O}.
\]
\begin{enumerate}[label=\textbf{\alph*)},leftmargin=2.2em,topsep=2pt,itemsep=2pt]
  \item Écris l'équation \textbf{ionique} (en ne gardant que les espèces qui changent) et identifie les
        deux couples avec leur nombre d'électrons.
  \item Équilibre les deux demi-équations en milieu acide.
  \item Combine-les pour obtenir l'\textbf{équation ionique globale}.
  \item Réintroduis les ions \ce{K+} et \ce{Cl^-} pour retrouver l'\textbf{équation moléculaire} complète.
  \item On propose : \textbf{(A)} le coefficient de \ce{MnCl2} vaut 1 ; \textbf{(B)} 12 électrons sont
        échangés ; \textbf{(C)} $\NO(\ce{Mn})$ dans \ce{KMnO4} vaut $+8$ ; \textbf{(D)} il faut ajouter
        8 \ce{H2O} dans le membre de droite. Laquelle est \textbf{correcte} ?
\end{enumerate}
\end{exo}

\begin{exo}[dichromate et acide nitreux — d'après annales 2022]
On considère la réaction non pondérée :
\[
\ce{Cr2O7^{2-} + HNO2 + H3O+ -> Cr^{3+} + NO3^- + H2O}.
\]
\begin{enumerate}[label=\textbf{\alph*)},leftmargin=2.2em,topsep=2pt,itemsep=2pt]
  \item Détermine $\NO(\ce{N})$ dans \ce{HNO2} et dans \ce{NO3^-}. Identifie les deux couples et le
        nombre d'électrons de chacun.
  \item Équilibre les deux demi-équations en milieu acide (avec \ce{H+}).
  \item Combine-les en une équation globale écrite avec \ce{H+}.
  \item Récris-la avec \ce{H3O+} et donne la suite des coefficients
        $(\ce{Cr2O7^{2-}},\ \ce{HNO2},\ \ce{H3O+}) \to (\ce{Cr^{3+}},\ \ce{NO3^-},\ \ce{H2O})$.
\end{enumerate}
\end{exo}

\begin{exo}[équilibrage complet en milieu basique — d'après livre QCM 14]
On considère la réaction en \textbf{milieu basique} :
\[
\ce{H2O + CrO4^{2-} + Fe -> Cr2O3 + Fe2O3 + OH^-}.
\]
\begin{enumerate}[label=\textbf{\alph*)},leftmargin=2.2em,topsep=2pt,itemsep=2pt]
  \item Identifie les deux couples (\ce{CrO4^{2-}}/\ce{Cr2O3} et \ce{Fe2O3}/\ce{Fe}) et le nombre
        d'électrons de chacun (attention aux 2 atomes de Cr et de Fe).
  \item Équilibre la demi-équation de \textbf{réduction} en milieu basique.
  \item Équilibre la demi-équation d'\textbf{oxydation} en milieu basique.
  \item Combine-les, donne l'équation globale équilibrée, puis calcule la \textbf{somme de tous les
        coefficients}. Le fer joue-t-il le rôle d'oxydant ou de réducteur ?
\end{enumerate}
\end{exo}

\end{document}
